% Secao 0 - Sumario executivo (1 pag). Quem le so isto entende o achado.
% Headline number unico, repetido no abstract e na conclusao.

\section*{Sumário executivo}
\addcontentsline{toc}{section}{Sumário executivo}

\noindent\textbf{Tese.} Desertos hospitalares medidos em quilômetros e desertos
medidos pelo fluxo real de pacientes não são o mesmo mapa --- e o gap entre eles
identifica os municípios que uma política de acesso baseada em distância
sistematicamente não enxerga.

\medskip
\noindent\textbf{O que fazemos.} Para os 5.492 municípios brasileiros (2015--2022),
confrontamos o deserto de \emph{distância} (km ao hospital ativo mais próximo)
com o deserto de \emph{fluxo} (burden efetivo de viagem, ponderado por onde os
pacientes de fato internam, construído de 94 milhões de internações do SIH). O
exercício é descritivo; nenhuma reivindicação causal é feita.

\medskip
\noindent\textbf{Três números que o relatório sustenta.}
\begin{itemize}
\item As duas medidas correlacionam-se apenas moderadamente ($r=0{,}55$): a
  distância ao hospital mais próximo explica cerca de 30\% da variação no burden
  efetivo de fluxo. O paciente típico viaja 32~km \emph{além} do hospital mais
  próximo.
\item Dos municípios sinalizados como deserto por qualquer dos dois conceitos,
  60\% o são por apenas um deles. Os mapas discordam sobre a maioria dos
  desertos.
\item \textbf{595 municípios, onde vivem cerca de 9,2 milhões de pessoas, têm um
  hospital por perto no mapa mas são desertos de fluxo} --- seus pacientes
  viajam, em média efetiva, 91~km, e 90\% das internações deixam o município.
  São invisíveis a uma política de distância.
\item O deserto de fluxo tem dois lados. Os mesmos municípios que perdem
  pacientes têm uma vez e meia menos médicos por habitante e menos especialistas
  --- apesar de terem leito ---, e sangram médicos ano após ano para os centros
  que absorvem seus pacientes. A ausência do cuidado e a de quem o presta
  caminham juntas.
\end{itemize}

\medskip
\noindent\textbf{A figura-síntese.} A Figura~\ref{fig:divergencia} pinta o
desacordo: em vermelho, os desertos de fluxo ocultos à distância, no interior do
Nordeste e do Centro-Oeste; em azul, onde a distância exagera o isolamento, na
Amazônia de viagem fluvial. Os perfis confirmam que o ponto cego recai sobre
municípios mais pobres e dependentes do SUS, e que ter um hospital local não
equivale a ter o cuidado de que se precisa --- o lado da oferta confirma por quê:
esses municípios têm menos médicos, menos especialistas e perdem médicos para os
mesmos hubs (Figura~\ref{fig:redes_medicas}). A pergunta sobre o custo em saúde
do isolamento exige um desenho causal e fica reservada a trabalho relacionado;
este relatório entrega a medida e a geografia.
